\clearpage
\chapter{曲式入门}

\classquote{夫人之立言，因字而生句，积句而成章，积章而成篇。}{《文心雕龙·章句》}{字生句、句积章、章成篇，和动机生乐句、乐句积段落、段落成全曲是同一个道理。曲式不是套模板，而是安排“记住什么、期待什么、何时兑现”。}

\label{sec:form-basics}

\subsection*{学完这一章能做什么}
\begin{itemize}
  \item 把一首歌从“连续的几段音乐”拆成有功能的段落，而不是只凭感觉说“这里换了一段”。
  \item 用重复、变化、对比三种手法，把四小节乐句组织成八小节、十六小节乃至整曲结构。
  \item 认出 AABA 与主歌—副歌两种最常见的流行结构，并说清各自的强处与风险。
  \item 设计一个能落地的段落过渡，并在夜半的编曲区用片段复制、粘贴与段落带把结构排出来。
\end{itemize}

\subsection{曲式：把时间分成可记住的层次}

曲式（form）是一首作品在时间上的组织方式。它回答的不是“这段好不好听”，而是“哪一段在负责什么、下一段为什么接在这里”。同一串和弦，按四小节一段重复三次，和按“主歌—副歌—主歌—副歌”排列，听感完全不同；差别不在材料，而在材料被放在什么位置、承担什么功能。

曲式有层次。最小的单位是动机（motive），通常一到两小节；几个动机组成乐句（phrase），常见四到八小节；几个乐句组成段落（section），常见八到三十二小节；若干段落再组成整曲。上一章讲的是“把动机长成乐句”，这一章讲的是“把乐句放进段落，把段落排成整曲”。

层次不是按小节数硬切的。一个八小节段落也可能只包含一个长乐句；一个四小节段落也可能是两个两小节乐句的并置。判断段落时，先看它有没有自己的开头、发展、收束，再看它和前后段是否承担不同功能。段落名（主歌、副歌、桥段）是功能名，不是长度名。

\subsection{重复与变化：四种基本组织手法}

音乐靠重复建立坐标，靠变化制造方向。完全不重复，听众记不住；完全不变化，听众会走神。最常用的四种手法如下。

\begin{center}
\small
\begin{tabular}{p{0.18\linewidth}p{0.34\linewidth}p{0.34\linewidth}}
\toprule
手法 & 做法 & 听感与风险 \\
\midrule
原样重复 & 同一乐句或段落再来一次，音高、节奏、和弦基本不动 & 建立记忆点、确认结构；连续两次以上容易停滞 \\
变化重复 & 保留骨架，改尾音、节奏、音区、配器或歌词 & 既熟悉又有推进；改得太多会失去“同一段”的身份 \\
模进/移位 & 整段材料按音级或半音整体挪高、挪低 & 产生方向感和推进感；连续模进容易变成机械爬升 \\
对比 & 换材料、换调、换节奏密度、换音区或换音色 & 划出段落边界、制造转折；对比过猛会像另一首歌 \\
\bottomrule
\end{tabular}

表 8-1：四种基本组织手法。先决定“这一段要让人记住什么”，再决定哪里重复、哪里变化。
\end{center}

一个实用的比例是“二比一”：两次可辨认的重复，配一次明确的变化。AABA 里的 A 出现三次，但第二次和第三次通常不是逐字重复；主歌—副歌里副歌反复出现，但每遍的歌词、配器或结尾可以不同。重复给耳朵地图，变化给耳朵路标。

\subsection{常见结构一：AABA}

AABA 是最经典的短篇结构之一，常见于三十二小节的歌曲：四个八小节段落，按 A—A—B—A 排列。A 段陈述主题，第二遍 A 通常换歌词或换配器，B 段（bridge）转去一个新方向，最后 A 回来收束。

\begin{center}
\begin{minipage}{0.94\linewidth}
\centering
\begin{tikzpicture}[x=1cm,y=1cm,font=\scriptsize]
  \foreach \i/\lab/\x/\a/\b in {0/A/0/1/8,1/A/3.3/9/16,2/B/6.6/17/24,3/A/9.9/25/32} {
    \draw[rounded corners=2pt,fill=gray!12] (\x,0) rectangle ++(3.0,0.9);
    \node at (\x+1.5,0.45) {\textbf{\lab}};
    \node[below,font=\tiny] at (\x+1.5,-0.04) {第 \a--\b 小节};
  }
  \draw[<->,gray] (0,1.05) -- (3.0,1.05) node[midway,above,font=\tiny]{陈述};
  \draw[<->,gray] (3.3,1.05) -- (6.3,1.05) node[midway,above,font=\tiny]{确认};
  \draw[<->,gray] (6.6,1.05) -- (9.6,1.05) node[midway,above,font=\tiny]{离开};
  \draw[<->,gray] (9.9,1.05) -- (12.9,1.05) node[midway,above,font=\tiny]{回来};
\end{tikzpicture}

图 8-1：三十二小节 AABA 的段落地图。四个八小节块分别承担陈述、确认、离开、回来；B 段不必是新调，但必须有可辨认的不同方向。
\end{minipage}
\end{center}

AABA 的强处是“记忆负担小、回来得干净”。听众只需要认住一个 A，就能在最后得到熟悉的收束；B 段提供了一个受控的偏离，不会把整首歌拆散。风险是平：如果三遍 A 的歌词、配器、节奏完全一样，最后一遍会像补时长；如果 B 段只是把 A 换个和弦继续重复，又缺少真正的“离开”。

写 AABA 时，先定 A 段的四小节问句与四小节答句，再把答句的尾音或节奏改一次，得到 A'；B 段可以转到关系小调、下属方向，或只换节奏密度与音区；最后一遍 A 不必原样返回，可以把结尾抬高半句、换配器或加和声，让“回来”比第一遍更有重量。

\subsection{常见结构二：主歌—副歌}

主歌—副歌（verse--chorus）是现代流行、摇滚与民谣最常用的结构。它不靠一个大段落反复，而靠两种功能交替：主歌负责叙事、铺垫与降低能量；副歌负责标题、情绪高点与记忆点。二者之间常有一个预副歌（pre-chorus）做能量爬坡，后面可以接桥段（bridge）或间奏。

一个常见的八段骨架是：前奏（intro）—主歌 1—副歌 1—主歌 2—副歌 2—桥段—副歌 3—尾奏（outro）。主歌每遍换歌词，副歌每遍基本保留同一句核心歌词；桥段离开主歌与副歌的和声/节奏惯性，最后一遍副歌把能量推到最高，再用短尾奏落地。

主歌与副歌的差别不只是“歌词位置”。主歌通常音区较低、节奏较说话、配器较薄；副歌音区更高、元音更开、鼓组与低频更满，标题句落在强拍。预副歌的功能是把这两端接起来：它往往缩短乐句、提高和声张力、加入上升音型或鼓点加速，让副歌的“开”有准备。

\begin{center}
\small
\begin{tabular}{p{0.16\linewidth}p{0.24\linewidth}p{0.22\linewidth}p{0.20\linewidth}}
\toprule
结构 & 强处 & 风险 & 适合 \\
\midrule
AABA & 记忆负担小，收束干净，适合短篇 & 三遍 A 容易平；B 段容易只是“换和弦” & 爵士标准曲、老式流行、短篇叙事 \\
主歌—副歌 & 情绪弧线清楚，副歌可反复强化 & 主歌容易写成“等副歌的过场”；桥段容易多余 & 现代流行、摇滚、民谣、短视频歌曲 \\
\bottomrule
\end{tabular}

表 8-2：AABA 与主歌—副歌的分工。结构没有高下，只有“这首歌想把重量放在哪里”。
\end{center}

写主歌—副歌时，先把副歌的标题句与和弦落点定死，再写主歌如何“不抢”副歌。主歌如果也用副歌的音区、密度和重音，副歌就没有地方可升；反过来，副歌如果只是主歌的加速版，也会缺少真正的情绪换挡。预副歌不是必需段，但当主歌与副歌的音区、节奏差得太远时，它是最便宜的桥。


\begin{center}
\begin{minipage}{0.94\linewidth}
\centering
\begin{tikzpicture}[x=1.35cm,y=0.85cm,font=\scriptsize]
  \draw[->,gray] (0,0) -- (7.6,0) node[right,font=\tiny]{段落};
  \draw[->,gray] (0,0) -- (0,4.2) node[above,font=\tiny]{能量};
  \foreach \x/\lab in {0/前奏,1/主1,2/副1,3/主2,4/副2,5/桥,6/副3,7/尾} {
    \draw[gray!30] (\x,0) -- (\x,3.8);
    \node[below,font=\tiny,rotate=30,anchor=east] at (\x,0) {\lab};
  }
  \draw[very thick] plot coordinates {(0,0.8) (1,1.5) (2,3.2) (3,1.8) (4,3.4) (5,1.2) (6,4.0) (7,0.5)};
  \foreach \x/\y in {0/0.8,1/1.5,2/3.2,3/1.8,4/3.4,5/1.2,6/4.0,7/0.5} {
    \fill (\x,\y) circle (1.5pt);
  }
\end{tikzpicture}

图 8-2：主歌—副歌的能量曲线。桥段主动下沉，最后一遍副歌再抬到最高，尾奏快速回落；曲线只表示相对能量，不是精确音高或音量。
\end{minipage}
\end{center}

\subsection{段落过渡：让下一段有理由开始}

段落之间最怕“硬切到另一首歌”。过渡不需要很长，但需要一个可听出来的信号：速度变化、和声张力、鼓点填空、音区上升、动机转接，或者一个明确的休止。下面四种是流行写作里最常用的。

\begin{center}
\small
\begin{tabular}{p{0.20\linewidth}p{0.38\linewidth}p{0.30\linewidth}}
\toprule
过渡 & 做法 & 注意 \\
\midrule
硬切 & 上一段结束，下一段直接在强拍开始 & 适合短促、有力的副歌；前后调性/音区差太大时会突兀 \\
鼓点填空 & 在最后一拍到两拍加军鼓、通鼓或加速单击 & 最通用的能量爬坡；填得太满会抢副歌第一拍 \\
上升/下降 & 用 riser、反向镲、噪音或弦乐滑音把能量推上去/放下来 & 上升音型要与副歌落点对齐；过长会变成等待 \\
共同和弦/动机桥 & 用两段都有的和弦、动机或旋律音做枢轴 & 最自然，也最容易被忽略；需要提前一拍埋下 \\
\bottomrule
\end{tabular}

表 8-3：四种常用段落过渡。先决定“下一段要从哪个拍子开始”，再倒推过渡的长度。
\end{center}

一个可操作的检查清单：过渡的最后两拍有没有明确的加速或停顿？下一段第一拍的重音是否落在听众预期的地方？过渡里使用的动机是否在下一段被接住？如果三个答案都是“没有”，这段过渡大概率只是把两段剪在一起。

\subsection{在夜半里怎么做}

夜半的编曲区适合做“结构排版”，不适合替你判断结构好坏。先在一张纸上写出段落与小节范围，再进工程摆片段。

\begin{enumerate}
  \item 先写结构表：例如 A 段第 1--8 小节、A' 段第 9--16 小节、B 段第 17--24 小节、A'' 段第 25--32 小节。把每段的功能（陈述/确认/离开/回来）写在后面。
  \item 选中一个片段，按 \texttt{Ctrl+C} 复制、\texttt{Ctrl+V} 粘贴；粘贴的落点是当前选中轨道的播放头位置，多个片段会保留彼此的相对间距。\texttt{Ctrl+X} 是剪切。
  \item 如果只想把选中的片段原样接在后面，按 \texttt{Ctrl+D}：没有音符选区时，它会复制编曲区选中的片段，副本从原片段结束处开始，并在片段名后加“ (Copy)”。注意：如果钢琴卷帘里还有音符选区，\texttt{Ctrl+D} 优先复制音符，不会复制整段。
  \item 用轨道头或混音台确认每条轨道的角色（鼓、贝斯、和弦、旋律）。当前界面的轨道名与片段名是只读的：轨道头/混音台没有改名输入框，片段名由新建逻辑生成（\texttt{Clip}、\texttt{Clip 2}）。需要改名时，走 MCP 的 \texttt{setTrackName} / \texttt{setClipName}；GUI 改名入口尚未提供。
  \item 编曲区顶部的段落带会显示工程里已有的 \texttt{sections}（名称、颜色、起止位置）；点击段落卡片可把播放头定位到该段起点。当前 GUI 不能在段落带上新建、拖动或改名，段落数据要由工程文件或 MCP 的 \texttt{setSection} 提供。
  \item 打开总览条，确认段落分界线与小节号对得上；如果段落带为空，说明这个工程还没有 \texttt{sections}，不要把它当成“界面坏了”。
\end{enumerate}

\screenshot{arrangement-aaba-structure}{编曲区 AABA 段落结构}

\textbf{范围说明：}夜半当前没有自动曲式分析、段落识别、段落自动命名，也没有“按结构生成整曲”的功能；段落带只读取工程里已有的 \texttt{sections}，不会从音频或 MIDI 自动推断。片段复制只复制编排位置与片段引用，不会自动改写歌词、配器或和弦；结构判断仍要靠耳朵与手动排版。

\subsection{练习}

\begin{enumerate}
  \item 把三十二小节的 AABA 写成小节范围，并指出 B 段承担什么功能。
  \par\noindent 答案：A 第 1--8 小节，A' 第 9--16 小节，B 第 17--24 小节，A'' 第 25--32 小节。B 段负责离开主方向，提供对比与转折；最后 A 回来收束。
  \item 一个八小节主歌与八小节副歌之间，你会把预副歌放在哪里？它需要完成什么？
  \par\noindent 答案：放在主歌结束到副歌开始之间，通常四到八小节；它要缩短乐句、提高和声张力或加速鼓点，把主歌的较低能量接到副歌的高点。
  \item 把一个四小节乐句原样重复四次，和把它做成 AABA，听感差别在哪里？
  \par\noindent 答案：原样四次只有确认、没有离开，容易平；AABA 在第三次换成 B 段，给耳朵一个受控的偏离，最后再回来。
  \item 为主歌到副歌设计一个两小节的过渡：写出你用的信号、落点和风险。
  \par\noindent 答案：不唯一。合格答案要有一个可听信号（鼓点填空、riser、共同和弦或动机桥）、明确的副歌第一拍落点，以及一个风险控制（例如 riser 不要盖过副歌第一拍）。
  \item 在夜半把第 1--4 小节的片段复制到第 5--8 小节，再复制到另一条轨道；写出你的操作和前置条件。
  \par\noindent 答案：选中片段后 \texttt{Ctrl+C}，选中目标轨道、把播放头放到第 5 小节，\texttt{Ctrl+V}；复制到另一条轨道时先切换选中轨道再粘贴。只想原样接在后面也可以用 \texttt{Ctrl+D}，但要先清掉钢琴卷帘里的音符选区。
\end{enumerate}

\subsection{帮助主题映射}

\begin{center}
\small
\begin{tabular}{p{0.26\linewidth}p{0.45\linewidth}p{0.18\linewidth}}
\toprule
本章小节 & 应用内帮助锚点 & 备注 \\
\midrule
曲式与层次 & \texttt{music-theory-form} & 概念 \\
重复与变化 & \texttt{music-theory-repetition-variation} & 表 8-1 \\
AABA & \texttt{music-theory-aaba} & 图 8-1 \\
主歌—副歌 & \texttt{music-theory-verse-chorus} & 表 8-2 \\
段落过渡 & \texttt{music-theory-transition} & 表 8-3 \\
在夜半里怎么做 & \texttt{music-theory-arrangement-structure} & 片段与段落带 \\
练习 & \texttt{music-theory-ch8-practice} & 答案在章末 \\
\bottomrule
\end{tabular}
\end{center}

\subsection*{延伸阅读}
\begin{itemize}
  \item \textbf{Open Music Theory, “The small binary.”} \url{https://openmusictheory.github.io/smallBinary.html} —— 用两段式讲重复、对比与回归，补本章最基础的结构骨架。
  \item \textbf{YouTube, “The Hidden Structure Behind Every Song You've Ever Heard.”} \url{https://www.youtube.com/watch?v=v8fXuLi9Rb4} —— 两分钟把音符、节奏、旋律、和声与结构的关系串起来，适合先看全景。
\end{itemize}

